% !PRINTED 66
\weberpage{66}
%<<CONT

\(1,2,\ldots,n\) in bestimmter Weise \(n\) reelle Zahlwerthe \(a_1,a_2,\ldots,a_n\) zuordnen, und aus diesen Zahlwerthen die \(\frac{n(n-1)}2\) Differenzen
\[
a_1-a_2,\quad a_1-a_3,\ldots
\]
bilden, wobei der niedrigere Index dem Minuenden angehören soll. Das Differenzenproduct
\begin{equation}
P=(a_1-a_2)(a_1-a_3)\cdots(a_1-a_n)
(a_2-a_3)\cdots(a_2-a_n)\cdots(a_{n-1}-a_n)
\tag{1}
\end{equation}
wird, wenn die \(a_1,a_2,\ldots,a_n\) von einander verschiedene Zahlwerthe sind, einen von Null verschiedenen Werth haben, z.~B. einen positiven, wenn \(a_1>a_2>a_3>\cdots>a_n\) angenommen war.

Wenn wir nun die Indices \(1,2,3,\ldots,n\) irgendwie unter einander vertauschen, also etwa von \(\mathfrak A\) zu \(\mathfrak A'\) übergehen, so geht \(P\) in
\begin{equation}
\begin{aligned}
P'={}&(a_{\alpha_1}-a_{\alpha_2})(a_{\alpha_1}-a_{\alpha_3})\cdots
(a_{\alpha_1}-a_{\alpha_n})\\
&\cdot(a_{\alpha_2}-a_{\alpha_3})\cdots(a_{\alpha_2}-a_{\alpha_n})
\cdots(a_{\alpha_{n-1}}-a_{\alpha_n})
\end{aligned}
\tag{2}
\end{equation}
über, und dies Product besteht, abgesehen vom Vorzeichen, aus denselben Factoren wie \(P\), d.~h. es ist \(P'\) entweder gleich \(P\) oder entgegengesetzt zu \(P\).

\gesperrt{I. Wir rechnen nun die Anordnung \(\mathfrak A'\) und also auch die Permutation, die \(\mathfrak A\) in \(\mathfrak A'\) verwandelt, zur ersten oder zur zweiten Art, je nachdem \(P\) mit \(P'\) gleich oder entgegengesetzt ist, so dass \(\mathfrak A\) selbst zur ersten Art gehört.}

\gesperrt{II. Durch eine einfache Transposition \((h,k)\), worin \(h,k\) irgend zwei der Ziffern \(1,2,\ldots,n\) bezeichnen, ändert sowohl \(P\) als \(P'\) sein Vorzeichen.}

Denn die Factoren, die \(h\) und \(k\) gar nicht enthalten, werden durch diese Transposition nicht berührt; dann haben wir in \(P\) und \(P'\) den Factor \(\pm(a_h-a_k)\) und die Factorenpaare
\[
\pm(a_h-a_\nu)(a_k-a_\nu),
\]
wo \(\nu\) die Reihe der Zahlen \(1,2,\ldots,n\), mit Ausnahme von \(h,k\), durchläuft. Der erstere Factor ändert aber sein Zeichen, während das Factorenpaar ungeändert bleibt bei der Transposition \((h,k)\). Daraus folgt:
%CONT>>
